\section{Discussion}
CertiPatch is designed for scope-bounded, replayable empirical guarantees rather than formal verification.
The protocol can fail when the patch family cannot express a feasible repair under tight collateral constraints; such cases must be reported fail-closed.
Scaling to larger models is primarily an engineering problem under the adapter contract and does not change the certificate semantics.
This submission is single-seed and compute-aware; we therefore do not claim robustness across seeds.
Baseline coverage is partial in some settings, so comparative claims are limited to completed matrix cells.
For coverage-bounded domains, certificates are explicitly bounded by the recorded coverage plan and do not imply out-of-scope correctness.

\paragraph{Interpretation of collateral.}
Collateral is highly task-dependent in this run profile: PARITY-4D achieves near-zero KL and low drift, while COMPARE-2D and BALANCE-PAREN-14 show larger measured drift under the same protocol.
COMPARE-6D-STRAT likewise shows high collateral movement under its coverage-bounded regime.
This is the main paper-risk axis and is explicitly scoped in our claims.
We do not claim a uniform bounded-collateral guarantee across specs.
This suggests that patch complexity alone is insufficient as a quality proxy; certification reports should always include both feasibility and collateral diagnostics.

\paragraph{Practical implication.}
The strongest claim supported here is operational: practitioners can replay, verify, and audit what was repaired and under which scope assumptions.
This is a stronger safety posture than reporting isolated task accuracy without certificate semantics.
